\subsection{AUTOMORPHISMS OF A SPLIT SEMI-SIMPLE LIE ALGEBRA}

Let E be a commutative group, and \(A = \bigoplus_{\gamma \in E} A^{\gamma}\) an E-graded algebra. For any homomorphism \(\varphi\) from the group E to the multiplicative group \(k^{*}\) , let \(f(\varphi)\) be the k-linear map from A to A whose restriction to each \(A^{\gamma}\) is the homothety with ratio \(\varphi(\gamma)\) ; it is clear that \(f(\varphi)\) is an automorphism of the graded algebra A, and that f is a homomorphism from the group \(\operatorname{Hom}(E, k^{*})\) to the group of automorphisms of the graded algebra A.

Let h be a splitting Cartan subalgebra of g, and R the root system of \((\mathfrak{g}, \mathfrak{h})\) . Recall that P(R) (resp. Q(R)) denotes the group of weights (resp. radical weights) of R. Put

\[
\mathrm{T} _ {\mathrm{P}} = \operatorname{Hom} (\mathrm{P} (\mathrm{R}), k ^ {*}) \quad \mathrm{T} _ {\mathrm{Q}} = \operatorname{Hom} (\mathrm{Q} (\mathrm{R}), k ^ {*}).
\]

We can consider \(g = g^{0} + \sum_{\alpha \in R} g^{\alpha}\) as a \(Q(R)\) -graded algebra. The preceding remarks define a canonical homomorphism from \(T_{Q}\) to \(\operatorname{Aut}(\mathfrak{g}, \mathfrak{h})\) , which will be denoted by f in this paragraph. In the other hand, the canonical injection from \(Q(R)\) to \(P(R)\) defines a homomorphism from \(T_{P}\) to \(T_{Q}\) , which will be denoted by q:

\[
\mathrm{T} _ {\mathrm{P}} \quad \stackrel {{q}} {{\longrightarrow}} \quad \mathrm{T} _ {\mathrm{Q}} \quad \stackrel {{f}} {{\longrightarrow}} \quad \operatorname{Aut} (\mathfrak {g}, \mathfrak {h}).
\]

If \(s \in \operatorname{Aut}(\mathfrak{g}, \mathfrak{h})\) , let \(s^*\) be the restriction of \({}^t(s|_{\mathfrak{h}})^{-1}\) to \(\mathrm{Q}(\mathrm{R})\) . Then, for all \(\varphi \in \mathrm{T}_{\mathrm{Q}}\) ,

\[
f (\varphi \circ s ^ {*}) = s ^ {- 1} \circ f (\varphi) \circ s.\tag{1}
\]

Indeed, let \(\gamma \in \mathrm{Q}(\mathrm{R})\) and \(x\in \mathfrak{g}^{\gamma}\) ; then \(sx\in \mathfrak{g}^{s^{*}\gamma}\) and

\[
f (\varphi \circ s ^ {*}) x = (\varphi \circ s ^ {*}) (\gamma). x = s ^ {- 1} (\varphi (s ^ {*} \gamma) s x) = (s ^ {- 1} \circ f (\varphi) \circ s) (x).
\]

PROPOSITION 2. The sequence of homomorphisms

\[
1 \longrightarrow \mathrm {T_ {Q}} \stackrel {f} {\longrightarrow} \mathrm{Aut} (\mathfrak {g}, \mathfrak {h}) \stackrel {\varepsilon} {\longrightarrow} \mathrm{A(R)} \longrightarrow 1
\]

is exact.

\begin{enumerate}
\def\labelenumi{\alph{enumi})}
\item
  Let \(\varphi \in \operatorname{Ker} f\) . Then \(\varphi(\alpha) = 1\) for all \(\alpha \in \mathrm{R}\) . Since \(\mathrm{R}\) generates the group \(\mathrm{Q}(\mathrm{R})\) , \(\varphi\) is the identity element of \(\mathrm{T}_{\mathrm{Q}}\) .
\item
  Let \(\varphi \in \mathrm{T}_{\mathrm{Q}}\) . The restriction of \(f(\varphi)\) to \(\mathfrak{h} = \mathfrak{g}^0\) is the identity, so
\end{enumerate}

\(\operatorname{Im} f \subset \operatorname{Ker} \varepsilon\) .

\begin{enumerate}
\def\labelenumi{\alph{enumi})}
\setcounter{enumi}{2}
\tightlist
\item
  Let \(s \in \operatorname{Ker} \varepsilon\) . Then \(s|_{\mathfrak{h}} = \mathrm{Id}_{\mathfrak{h}}\) . For all \(\alpha \in R\) , we have \(s(\mathfrak{g}^{\alpha}) = \mathfrak{g}^{\alpha}\) , and there exists a \(t_{\alpha} \in k^{*}\) such that \(sx = t_{\alpha}x\) for all \(x \in g^{\alpha}\) . Writing down the condition that \(s \in \operatorname{Aut}(\mathfrak{g})\) , we obtain the relations
\end{enumerate}

\[
t _ {\alpha} t _ {- \alpha} = 1 \qquad \text{for all } \alpha \in \mathrm{R}
\]

\[
t _ {\alpha} t _ {\beta} = t _ {\alpha + \beta} \quad \text { when } \alpha , \beta , \alpha + \beta \in \mathrm{R}.
\]

Under these conditions, there exists \(\varphi \in \mathrm{T}_{\mathrm{Q}}\) such that \(\varphi (\alpha) = t_{\alpha}\) for all \(\alpha \in \mathrm{R}\) (Chap. VI, §1, no. 6, Cor. 2 of Prop. 19). Then \(s = f(\varphi)\) . Hence, \(\operatorname{Ker} \varepsilon \subset \operatorname{Im} f\) .

\begin{enumerate}
\def\labelenumi{\alph{enumi})}
\setcounter{enumi}{3}
\tightlist
\item
  The image of \(\operatorname{Aut}(\mathfrak{g},\mathfrak{h})\) under \(\varepsilon\) contains \(\mathrm{W}(\mathrm{R})\) by §2, no. 2, Cor. of Th. 2, and contains \(\operatorname{Aut}(\mathrm{R},\mathrm{B})\) by Prop. 1. Hence this image is equal to \(\mathrm{A}(\mathrm{R})\) .
\end{enumerate}

COROLLARY 1. Let \((\mathrm{B}, (X_{\alpha})_{\alpha \in \mathrm{B}})\) be a framing of \((\mathfrak{g}, \mathfrak{h})\) . Let \(\mathrm{G}\) be the set of \(s \in \operatorname{Aut}(\mathfrak{g}, \mathfrak{h})\) that leave the framing invariant. Then \(\operatorname{Aut}(\mathfrak{g}, \mathfrak{h})\) is the semi-direct product of \(\mathrm{G}\) and \(\varepsilon^{-1}(\mathrm{W}(\mathrm{R}))\) .

Indeed, \(\mathrm{G} \cap \varepsilon^{-1}(\mathrm{W}(\mathrm{R})) = \{1\}\) by Prop. 1, and

\[
\operatorname{Aut} (\mathfrak {g}, \mathfrak {h}) = \mathrm{G}. \varepsilon^ {- 1} (\mathrm{W(R)})
\]

since \(\varepsilon\) is surjective (Prop. 2).

COROLLARY 2. The group \(\varepsilon^{-1}(\mathrm{W}(\mathrm{R}))\) operates simply-transitively on the set of framings of \((\mathfrak{g},\mathfrak{h})\) .

Indeed, \(\operatorname{Aut}(\mathfrak{g},\mathfrak{h})\) operates transitively on the set of framings of \((\mathfrak{g},\mathfrak{h})\) by §4, no. 4, Th. 2. Cor. 2 now follows from Cor. 1.

COROLLARY 3. Let B be a basis of R. The group \(\operatorname{Ker} \varepsilon = f(\mathrm{T}_{\mathrm{Q}})\) operates simply-transitively on the set of framings of \((\mathfrak{g}, \mathfrak{h})\) of the form \((\mathrm{B}, (X_{\alpha})_{\alpha \in \mathrm{B}})\) .

This follows immediately from Prop. 2.

Let \(\alpha \in \mathrm{R}\) , \(X_{\alpha} \in \mathfrak{g}^{\alpha}\) , \(X_{-\alpha} \in \mathfrak{g}^{-\alpha}\) be such that \([X_{\alpha}, X_{-\alpha}] = -H_{\alpha}\) . We have seen (§2, no. 2, Th. 2) that, for all \(t \in k^{*}\) , the restriction of the elementary automorphism

\[
\theta_ {\alpha} (t) = e ^ {\mathrm{ad} t X _ {\alpha}} e ^ {\mathrm{ad} t ^ {- 1} X _ {- \alpha}} e ^ {\mathrm{ad} t X _ {\alpha}}
\]

to \(\mathfrak{h}\) is the transpose of \(s_{\alpha}\) ; so \(\varepsilon(\theta_{\alpha}(t)) = s_{\alpha}\) and consequently \(\theta_{\alpha}(t)\theta_{\alpha}(-1) \in \operatorname{Ker} \varepsilon\) .

Lemma 1. Let \(\alpha \in \mathrm{R}\) and \(t\in k^{*}\) . Let \(\varphi\) be the homomorphism \(\lambda \mapsto t^{\lambda (H_{\alpha})}\) from \(\mathrm{Q(R)}\) to \(k^*\) . Then \(f(\varphi) = \theta_{\alpha}(t)\theta_{\alpha}(-1)\) .

Let \(\rho\) be the representation of \(\mathfrak{sl}(2,k)\) on \(\mathfrak{g}\) associated to \(X_{\alpha}\) . Let \(\pi\) be the representation of \(\mathbf{SL}(2,k)\) compatible with \(\rho\) . Introduce the notations \(\theta(t), h(t)\) of §1, no. 5. Since \(\rho(H) = \operatorname{ad} H_{\alpha}\) , the elements of \(\mathfrak{g}^{\lambda}\) are of weight \(\lambda(H_{\alpha})\) for \(\rho\) . By §2, no. 2, \(\theta_{\alpha}(t)\theta_{\alpha}(-1) = \pi(\theta(t)\theta(-1)) = \pi(h(t))\) . Hence the restriction of \(\theta_{\alpha}(t)\theta_{\alpha}(-1)\) to \(\mathfrak{g}^{\lambda}\) is the homothety of ratio \(t^{\lambda(H_{\alpha})}\) (§1, no. 5, Prop. 6), hence the lemma.

PROPOSITION 3. The image of the composite homomorphism

\[
\mathrm{T} _ {\mathrm{P}} \quad \xrightarrow {q} \quad \mathrm{T} _ {\mathrm{Q}} \quad \xrightarrow {f} \quad \operatorname{Aut} (\mathfrak {g}, \mathfrak {h})
\]

is contained in \(\mathrm{Aut}_e(\mathfrak{g})\) .

Let B be a basis of R. Then \((H_{\alpha})_{\alpha\in\mathrm{B}}\) is a basis of \(R^{\vee}\) , and the dual basis of \((H_{\alpha})_{\alpha\in\mathrm{B}}\) in \(h^{*}\) is a basis of the group P(R). Hence the group \(T_{P}\) is generated by the homomorphisms \(\lambda\mapsto t^{\lambda(H_{\alpha})}\;(t\in k^{*},\alpha\in\mathrm{B})\) . If \(\varphi\) is the restriction of such a homomorphism to Q(R), Lemma 1 proves that \(f(\varphi)\in\operatorname{Aut}_{e}(\mathfrak{g})\) , hence the proposition.

Let \(\bar{k}\) be an algebraic closure of \(k\) . The map which associates to any automorphism \(s\) of \(\mathfrak{g}\) the automorphism \(s \otimes 1\) of \(\mathfrak{g} \otimes_k \bar{k}\) is an injective homomorphism from \(\mathrm{Aut}(\mathfrak{g})\) to \(\mathrm{Aut}(\mathfrak{g} \otimes_k \bar{k})\) . We denote by \(\mathrm{Aut}_0(\mathfrak{g})\) the normal subgroup of \(\mathrm{Aut}(\mathfrak{g})\) which is the inverse image of \(\mathrm{Aut}_e(\mathfrak{g} \otimes_k \bar{k})\) under this homomorphism; this is the set of automorphisms of \(\mathfrak{g}\) that become elementary on extending the base field from \(k\) to \(\bar{k}\) . It is clear that \(\mathrm{Aut}_e(\mathfrak{g})\) is independent of the choice of \(\bar{k}\) , and that \(\mathrm{Aut}_e(\mathfrak{g}) \subset \mathrm{Aut}_0(\mathfrak{g})\) . The groups \(\mathrm{Aut}_0(\mathfrak{g})\) and \(\mathrm{Aut}_e(\mathfrak{g})\) can be distinct (Chap. VII, §13, no. 1). If \(\mathfrak{h}\) is a Cartan subalgebra of \(\mathfrak{g}\) , put

\[
\operatorname{Aut} _ {e} (\mathfrak {g}, \mathfrak {h}) = \operatorname{Aut} _ {e} (\mathfrak {g}) \cap \operatorname{Aut} (\mathfrak {g}, \mathfrak {h}), \quad \operatorname{Aut} _ {0} (\mathfrak {g}, \mathfrak {h}) = \operatorname{Aut} _ {0} (\mathfrak {g}) \cap \operatorname{Aut} (\mathfrak {g}, \mathfrak {h}).
\]

Lemma 2. Let \(\mathfrak{h}\) be a splitting Cartan subalgebra of \(\mathfrak{g}\) , and \(s \in \operatorname{Aut}_0(\mathfrak{g}, \mathfrak{h})\) . Assume that the restriction of \(s\) to \(\sum_{\alpha \in \mathrm{R}} \mathfrak{g}^\alpha\) does not have 1 as an eigenvalue. Then \(\varepsilon(s) = 1\) .

By extension of \(k\) , we are reduced to the case where \(s \in \mathrm{Aut}_e(\mathfrak{g}, \mathfrak{h})\) . The dimension of the nilspace of \(s - 1\) is at least \(\dim \mathfrak{h}\) (Chap. VII, §4, no. 4, Prop. 9). Hence \((s - 1)|_{\mathfrak{h}}\) is nilpotent. Since \(s|_{\mathfrak{h}} \in \mathrm{A}(\mathrm{R}^\vee)\) , \(s|_{\mathfrak{h}}\) is of finite order, and hence semi-simple (Chap. V, Appendix, Prop. 2). Consequently, \((s - 1)|_{\mathfrak{h}} = 0\) , which proves that \(\varepsilon(s) = 1\) .

Lemma 3. (i) Let \(m = (\mathrm{P}(\mathrm{R}) : \mathrm{Q}(\mathrm{R}))\) . If \(\varphi\) is the \(m\) th power of an element of \(\mathrm{T}_{\mathrm{Q}}\) , then \(\varphi \in q(\mathrm{T}_{\mathrm{P}})\) .

\begin{enumerate}
\def\labelenumi{(\roman{enumi})}
\setcounter{enumi}{1}
\tightlist
\item
  If \(k\) is algebraically closed, \(q(\mathrm{T}_{\mathrm{P}}) = \mathrm{T}_{\mathrm{Q}}\) .
\end{enumerate}

There exist a basis \((\lambda_{1},\ldots,\lambda_{l})\) of P(R) and integers \(n_{1}\geq1,\ldots,n_{l}\geq1\) such that \((n_{1}\lambda_{1},\ldots,n_{l}\lambda_{l})\) is a basis of Q(R). We have \(m=n_{1}\ldots n_{l}\) . Let \(\psi\in T_{Q}\) and put \(\psi(n_{1}\lambda_{1})=t_{1},\ldots,\psi(n_{l}\lambda_{l})=t_{l}\) . For \(i=1,\ldots,l\) , put \(m_{i}=\) \(\prod_{j \neq i} n_j\) . Let \(\chi\) be the element of \(T_P\) such that \(\chi(\lambda_1) = t_1^{m_1}, \ldots, \chi(\lambda_l) = t_l^{m_l}\) . Then

\[
\chi (n _ {i} \lambda_ {i}) = t _ {i} ^ {m _ {i} n _ {i}} = t _ {i} ^ {m} = (\psi^ {m}) (n _ {i} \lambda_ {i})
\]

so \(\chi|_{\mathrm{Q}(\mathrm{R})}=\psi^{m}\) . This proves (i). If k is algebraically closed, every element of \(k^{*}\) is the mth power of an element of \(k^{*}\) , so every element of \(T_{Q}\) is the mth power of an element of \(T_{Q}\) ; hence, (ii) follows from (i).

PROPOSITION 4. We have \(f(\mathrm{T}_{\mathrm{Q}}) \subset \mathrm{Aut}_0(\mathfrak{g}, \mathfrak{h})\) and \(\varepsilon^{-1}(\mathrm{W}(\mathrm{R})) = \mathrm{Aut}_0(\mathfrak{g}, \mathfrak{h})\) .

\begin{enumerate}
\def\labelenumi{\alph{enumi})}
\tightlist
\item
  Let \(\varphi \in \mathrm{T}_{\mathrm{Q}}\) and let \(\bar{k}\) be an algebraic closure of \(k\) . By Lemma 3, \(\varphi\) extends to an element of \(\operatorname{Hom}(\mathrm{P}(\mathrm{R}), k^*)\) . By Prop. 3,
\end{enumerate}

\[
f (\varphi) \otimes 1 \in \operatorname{Aut} _ {e} (\mathfrak {g} \otimes_ {k} \bar {k}, \mathfrak {h} \otimes_ {k} \bar {k}).
\]

Hence \(f(\varphi) \in \mathrm{Aut}_0(\mathfrak{g},\mathfrak{h})\) , and \(\operatorname{Ker} \varepsilon \subset \mathrm{Aut}_0(\mathfrak{g},\mathfrak{h})\) .

\begin{enumerate}
\def\labelenumi{\alph{enumi})}
\setcounter{enumi}{1}
\item
  The image of \(\mathrm{Aut}_e(\mathfrak{g},\mathfrak{h})\) under \(\varepsilon\) contains \(\mathrm{W(R)}\) (§2, no. 2, Cor. of Th. 2). In view of a), we see that \(\varepsilon^{-1}(\mathrm{W(R)})\subset \mathrm{Aut}_0(\mathfrak{g},\mathfrak{h})\) .
\item
  It remains to prove that \(\mathrm{Aut}_0(\mathfrak{g},\mathfrak{h})\subset \varepsilon^{-1}(\mathrm{W}(\mathrm{R}))\) . In view of b), it suffices to prove that \(\varepsilon (\mathrm{Aut}_0(\mathfrak{g},\mathfrak{h}))\cap \mathrm{Aut}(\mathrm{R},\mathrm{B})\) , where B denotes a basis of R, reduces to \{1\}.
\end{enumerate}

Let \(s \in \mathrm{Aut}_0(\mathfrak{g}, \mathfrak{h})\) be such that \(\varepsilon(s) \in \mathrm{Aut}(\mathrm{R}, \mathrm{B})\) . The subgroup of \(\mathrm{A}(\mathrm{R})\) generated by \(\varepsilon(s)\) has a finite number of orbits on \(\mathrm{R}\) . Let \(\mathrm{U}\) be such an orbit, of cardinal \(r\) , and \(\mathfrak{g}^{\mathrm{U}} = \sum_{\beta \in \mathrm{U}} \mathfrak{g}^{\beta}\) . Let \(\beta_1 \in \mathrm{U}\) , and put \(\beta_i = \varepsilon(s)^{i-1}\beta_1\) for \(1 \leq i \leq r\) , so that \(\mathrm{U} = \{\beta_1, \ldots, \beta_r\}\) . Let \(X_{\beta_1}\) be a non-zero element of \(\mathfrak{g}^{\beta_1}\) , and put \(X_{\beta_i} = s^{i-1}X_{\beta_1}\) for \(1 \leq i \leq r\) . There exists \(c_{\mathrm{U}} \in k^*\) such that \(s^r X_{\beta_1} = c_{\mathrm{U}}X_{\beta_1}\) , hence \(s^r X_{\beta_i} = c_{\mathrm{U}}X_{\beta_i}\) for all \(i\) , and consequently \(s^r |_{\mathfrak{g}^{\mathrm{U}}} = c_{\mathrm{U}}.1\) . Let \(\varphi \in T_{\mathrm{Q}}\) , and \(s' = s \circ f(\varphi)\) , which by \(a)\) is an element of \(\mathrm{Aut}_0(\mathfrak{g}, \mathfrak{h})\) . We have \(s'^r |_{\mathfrak{g}^{\mathrm{U}}} = c_{\mathrm{U}}'.1\) , where

\[
c _ {\mathrm{U}} ^ {\prime} = c _ {\mathrm{U}} \prod_ {i = 1} ^ {r} \varphi (\beta_ {i}) = c _ {\mathrm{U}} \varphi \left(\sum_ {i = 1} ^ {r} \beta_ {i}\right).
\]

Put \(\mathrm{B} = \{\alpha_1, \ldots, \alpha_l\}\) and \(\sum_{i=1}^{r} \beta_i = \sum_{j=1}^{l} m_j^{\mathrm{U}} \alpha_j\) . Since \(\varepsilon(s) \in \mathrm{Aut}(\mathrm{R}, \mathrm{B})\) , the \(m_j^{\mathrm{U}}\) are integers of the same sign and not all zero. We have

\[
c _ {\mathrm{U}} ^ {\prime} = c _ {\mathrm{U}} \prod_ {j = 1} ^ {l} \varphi (\alpha_ {j}) ^ {m _ {j} ^ {\mathrm{U}}}.
\]

Now \(\varphi\) can be chosen so that \(c_{\mathrm{U}}^{\prime} \neq 1\) for every orbit U; indeed, this reduces to choosing elements \(\varphi(\alpha_1) = t_1, \ldots, \varphi(\alpha_l) = t_l\) of \(k^*\) which are not annihilated by a finite number of polynomials in \(t_1, \ldots, t_l\) , not identically zero. For such a choice of \(\varphi\) , \(\varepsilon(s') = 1\) by Lemma 2, so

\[
\varepsilon (s) = \varepsilon (s ^ {\prime}) \varepsilon (f (\varphi)) ^ {- 1} = 1.
\]

COROLLARY. Let B be a basis of R. The group \(\operatorname{Aut}(\mathfrak{g},\mathfrak{h})\) is isomorphic to the semi-direct product of the groups \(\operatorname{Aut}(\mathrm{R},\mathrm{B})\) and \(\operatorname{Aut}_{0}(\mathfrak{g},\mathfrak{h})\) .

This follows from Prop. 1, Cor. 1 of Prop. 2, and Prop. 4.

Remark. Let \(\varepsilon^{\prime},\varepsilon^{\prime\prime}\) be the restrictions of \(\varepsilon\) to \(\operatorname{Aut}_{0}(\mathfrak{g},\mathfrak{h}),\operatorname{Aut}_{e}(\mathfrak{g},\mathfrak{h})\) . Let \(f^{\prime}\) be the homomorphism from \(T_{P}\) to \(\operatorname{Aut}_{e}(\mathfrak{g},\mathfrak{h})\) induced by f via the canonical injection from \(Q(R)\) to \(P(R)\) . In the preceding we have established the following commutative diagram:

\[
\begin{array}{c c c c c c c c c} 1 & \longrightarrow & \mathrm{T} _ {\mathrm{Q}} & \stackrel {{f}} {{\longrightarrow}} & \mathrm{Aut} (\mathfrak {g}, \mathfrak {h}) & \stackrel {{\varepsilon}} {{\longrightarrow}} & \mathrm{A} (\mathrm{R}) & \longrightarrow & 1 \\ & & \uparrow & & \uparrow & & \uparrow & & \\ 1 & \longrightarrow & \mathrm{T} _ {\mathrm{Q}} & \stackrel {{f}} {{\longrightarrow}} & \mathrm{Aut} _ {0} (\mathfrak {g}, \mathfrak {h}) & \stackrel {{\varepsilon^ {\prime}}} {{\longrightarrow}} & \mathrm{W} (\mathrm{R}) & \longrightarrow & 1 \\ & & \uparrow_ {q} & & \uparrow & & \uparrow & & \\ & & \mathrm{T} _ {\mathrm{P}} & \stackrel {{f ^ {\prime}}} {{\longrightarrow}} & \mathrm{Aut} _ {e} (\mathfrak {g}, \mathfrak {h}) & \stackrel {{\varepsilon^ {\prime \prime}}} {{\longrightarrow}} & \mathrm{W} (\mathrm{R}) & \longrightarrow & 1 \end{array}
\]

in which the vertical arrows other than q denote the canonical injections. We have seen (Prop. 2 and 4) that the first two rows are exact. In the third row, the homomorphism \(\varepsilon''\) is surjective (§2, no. 2, Cor. of Th. 2); it can be shown that its kernel is \(f'(\mathrm{T}_{\mathrm{P}})\) (§7, Exerc. 26 d)).

